\subsection{Dual Attention: Modality and Position Axes}\label{sec:methodology-fusion}

How a model combines image and label evidence, and how context and target
interact, are two related design choices this section addresses together.
Medverse~\cite{huMedverseUniversalModel2025a} concatenates image and label
channels before a shared U-Net; context and target then interact through a
Blockwise Cross-Attention Module (BAM) placed at two points in that network
(once encoder-side, once decoder-side), which pools its queries and keys to
block resolution before attending -- trading positional precision for a
cheaper attention cost -- while its values stay at full resolution.
Iris~\cite{gaoShowSegmentUniversal2025} fuses image and label features
once, through a pixel shuffle--unshuffle operation, then lets context
inform the target through a single bidirectional cross-attention step
against a compressed, per-class task embedding computed once and reused
across every query. Both commit every design choice -- fusion point and
context--target interaction alike -- to a single, early step, computed
once. \autoref{fig:fusion-family} contrasts Medverse's dual,
separately-weighted per-row encoders against the shared-encoder family
both our baseline and final design belong to.

\begin{figure}[t]
  \centering
  \includegraphics[width=\linewidth]{arch_fusion_family.pdf}
  \caption{Two families of context/target coupling, shown for a single
  context volume ($K{=}1$) and 3 of its $N$ tokens per stream. Top:
  Medverse gives every row its own, separately-weighted convolutional
  block, re-run at every step and coupled only through cross-attention
  (blue fan). Bottom: the family both our baseline and final design
  belong to instead pushes every row through one shared encoder, run
  once, before any repeated cross-attention step --
  \autoref{fig:biaxial} details the two ways we combine image and label
  evidence within this shared-encoder family.}
  \label{fig:fusion-family}
\end{figure}
\todo{Draft diagram, ported from the conference draft; finalize once the layer factorization is settled.}

We instead keep image tokens $z^{\text{img}}$ and label tokens
$z^{\text{lbl}}$ (\autoref{sec:methodology-overview}) as separate streams
throughout the network and let them exchange information repeatedly,
through \emph{dual attention} -- following the general pattern of
Fu et al.~\cite{fuDualAttentionNetwork2019}, who fuse a position-wise and a
channel-wise attention branch at every layer of a scene-segmentation
network -- with two attention axes applied at every layer, on the full,
uncompressed token grid:

\begin{itemize}[leftmargin=1.4em,itemsep=0.1em]
  \item \textbf{Modality axis}: at each cell, the image and label token
    attend to each other, letting label evidence sharpen the image
    representation and vice versa, independently per cell.
  \item \textbf{Position axis}: every cell of every volume -- context,
    target, and the register tokens -- attends over the entire token
    sequence at once, using 3D axial rotary position embeddings so
    attention is a function of physical distance between cells rather
    than raw grid index. This single joint attention performs
    within-volume self-attention and cross-context matching together,
    since tokens from every volume already share one sequence.
\end{itemize}

\autoref{fig:biaxial} contrasts this dual attention against a
single-fusion-step baseline within the same shared-encoder family.

\begin{figure}[t]
  \centering
  \includegraphics[width=\linewidth]{arch_fusion_axis.pdf}
  \caption{Two variants within the shared-encoder family
  (\autoref{fig:fusion-family}), shown for a single context volume
  ($K{=}1$ for clarity) and 3 of its $N$ spatial tokens per stream. Top:
  early fusion (our baseline) -- image and label tokens are merged once
  ($\oplus$), then the fused token grid attends across cells and
  contexts for $L$ layers (the position axis). Bottom: dual attention
  (ours) -- image and label tokens instead attend to each other directly
  at every layer (the modality axis, blue double arrows), alongside the
  same cross-cell/cross-context attention (the position axis), repeated
  together for $L$ layers. In both, the target row (green image tokens)
  cross-attends read-only into context (pink tokens); its label-stream
  input is the query prior (\autoref{sec:methodology-cascade}) rather
  than ground truth, which is why its label token is gray rather than a
  role color.}
  \label{fig:biaxial}
\end{figure}
\todo{Draft diagram, ported from the conference draft; finalize once the layer factorization is settled.}

Only the target's post-attention \emph{image} token is read out for
decoding; the label token is discarded after the last layer, since it
started from a placeholder rather than real image evidence for the target
row. The decoded image tokens are projected back to the coarse token grid
and progressively fused with the encoder's fine-resolution skip features
(\autoref{sec:methodology-overview}) via upsampling and convolution, one step
per fine stage, into the final full-resolution mask logits.
